% Figure: Fermi-Dirac occupation f(E) at two temperatures against the Boltzmann
% approximation. At T=0 the occupation is a step at the Fermi level; at finite T
% it softens over a width of a few k_B T. The Boltzmann tail exp(-(E-E_F)/kT)
% tracks Fermi-Dirac closely once E is more than a few k_B T above E_F, which is
% the regime the diode and solar-cell algebra of the chapter assumes.
% Curves drawn from explicit sample points; no computed exponentials in TikZ.
\begin{figure}[htbp]
\centering
\begin{tikzpicture}[font=\scriptsize]
\begin{axis}[
    width=11cm, height=6.4cm,
    xlabel={energy relative to Fermi level $(E-E_F)/k_B T_{300}$},
    ylabel={occupation $f(E)$},
    xmin=-6, xmax=6,
    ymin=0, ymax=1.05,
    xtick={-6,-4,-2,0,2,4,6},
    ytick={0,0.2,0.4,0.6,0.8,1.0},
    grid=both,
    grid style={enggray!25},
    tick label style={font=\tiny},
    label style={font=\scriptsize},
    axis line style={enggray},
    legend style={font=\tiny, at={(0.98,0.98)}, anchor=north east,
      draw=enggray, fill=white},
    every axis plot/.append style={very thick},
]
% Fermi-Dirac at T=300 K (x in units of k_B T_300): f = 1/(1+exp(x))
\addplot[engblue] coordinates {
  (-6,0.9975) (-5,0.9933) (-4,0.9820) (-3,0.9526) (-2,0.8808)
  (-1,0.7311) (0,0.5000) (1,0.2689) (2,0.1192) (3,0.0474)
  (4,0.0180) (5,0.0067) (6,0.0025)
};
\addlegendentry{Fermi-Dirac, $300\,\si{\kelvin}$}
% Fermi-Dirac at a hotter T (1.8x): softer step. f=1/(1+exp(x/1.8))
\addplot[engred, dashed] coordinates {
  (-6,0.9644) (-5,0.9398) (-4,0.9010) (-3,0.8425) (-2,0.7549)
  (-1,0.6358) (0,0.5000) (1,0.3642) (2,0.2451) (3,0.1575)
  (4,0.0990) (5,0.0602) (6,0.0356)
};
\addlegendentry{Fermi-Dirac, $540\,\si{\kelvin}$}
% Boltzmann tail at 300 K: exp(-x), only meaningful for x>0, clipped at 1
\addplot[engorange, densely dotted] coordinates {
  (0,1.0000) (1,0.3679) (2,0.1353) (3,0.0498) (4,0.0183)
  (5,0.0067) (6,0.0025)
};
\addlegendentry{Boltzmann tail $e^{-x}$}
% Fermi level guide
\draw[enggray, dashed] (axis cs:0,0) -- (axis cs:0,1.0);
\node[enggray, anchor=south west, font=\tiny] at (axis cs:0.1,0.92)
  {$E=E_F$};
% region where Boltzmann is accurate
\node[engorange, anchor=west, font=\tiny] at (axis cs:2.5,0.55)
  {Boltzmann accurate ($E-E_F\gtrsim 3k_BT$)};
\end{axis}
\end{tikzpicture}
\caption[Fermi-Dirac occupation and the Boltzmann limit]{Probability that a
state of energy $E$ is occupied by an electron, the Fermi-Dirac function
$f(E)=1/[1+\exp((E-E_F)/k_BT)]$, shown at room temperature and at nearly twice
that. At absolute zero the occupation is a sharp step at the Fermi level; finite
temperature softens it over a width of a few $k_BT$, and the higher curve is the
broader. The dotted line is the Boltzmann approximation $\exp(-(E-E_F)/k_BT)$,
which drops the $+1$ in the denominator. The two track within a percent once the
energy lies more than about three $k_BT$ above the Fermi level, the
non-degenerate regime in which the diode law, the intrinsic-carrier formula, and
the solar-cell voltage of this chapter are written.}
\label{fig:vol04ch12:fermi-dirac}
\end{figure}
